   % Latex Beamer template following CERN template guidelines (or trying!)

\documentclass[aspectratio=169]{beamer}
\usepackage{xcolor}
\usepackage{graphicx}
\usepackage{multicol}
\usepackage{tikz}

% Code listings with syntax highlighting
%  Require Pygments
\usepackage{minted}

\usetheme{LISA}

%%%%%%%%%%%%%%%%%%%%%%
%  Author and title  %
%%%%%%%%%%%%%%%%%%%%%%


% Preamble
\title[]{\bfseries \color{white} \textit{Ab initio} my title here}
\subtitle{\color{white} LISA academic day}
\author[Author]{\color{white} J. S. Andrews}

% Uncomment this to define a presentation date distinct from \today
% \def\mydate{20 Feb 2000}


\usebackgroundtemplate{   
\includegraphics[width=\paperwidth,height=\paperheight]{logos/footer.png}
}
\begin{document}

{
\setbeamertemplate{footline}{}
\setbeamertemplate{headline}{}
\usebackgroundtemplate{\includegraphics[height=\paperheight,width=\paperwidth]{logos/lisa_background.jpg}}
\begin{frame}[plain]
  \titlepage
\end{frame}
}

%    \setcounter{framenumber}{0}

    %%%%%%%%%%%%%%%%%%%%%%%
    %  Table of contents  %
    %%%%%%%%%%%%%%%%%%%%%%%
    

    \frame{
        \frametitle{Table of Contents}
        \begin{multicols}{2}
            \tableofcontents
        \end{multicols}
    }


    %%%%%%%%%%%%%%%%%%
    %  Introduction  %
    %%%%%%%%%%%%%%%%%%
    

    \section{Introduction}

    \frame{
        \frametitle{This is the first slide}
        %Content goes here
        \begin{block}{Ideas}
            See how this loooks
        \end{block}
        \begin{itemize}
            \item Item 1
            \item Item 2
        \end{itemize}
    }

    % This slide has code block with syntax highlighting using the Minted package
    %  Minted requires Pygments and building the PDF with the --shell-escape argument
    %  for pdflatex
    \begin{frame}[fragile]
        \frametitle{This slide has code blocks}
		\begin{minted}[
			frame=lines,
			framesep=2mm,
			baselinestretch=1.2,
			fontsize=\footnotesize,
			linenos
			]{python}
        import numpy as np
         
        def incmatrix(genl1,genl2):
            m = len(genl1)
            n = len(genl2)
            M = None #to become the incidence matrix
            VT = np.zeros((n*m,1), int)  #dummy variable
        \end{minted}
    \end{frame}


    \subsection{Intro subsection}
    \frame{
        \frametitle{This is the second slide}
        \framesubtitle{A bit more information about this}
        %More content goes here
        \begin{alertblock}{Alert block}
            Alert text
        \end{alertblock}
    }
    
    \section{Section 1}

    \frame{
        \frametitle{Generic slide}
        \framesubtitle{A bit more information about this}
        %More content goes here
    }
    \section{Section 2}

    \frame{
        \frametitle{Generic slide}
        \framesubtitle{A bit more information about this}
        %More content goes here
    }

    
    \section{Conclusion}
    
    \frame{
        \frametitle{Generic slide}
        \framesubtitle{A bit more information about this}
        %More content goes here
    }
    
    \section{Acknowledgements}
    
    \frame{
        \frametitle{Generic slide}
        \framesubtitle{A bit more information about this}
        %More content goes here
    }

    \section{}
    \lisaSplashWhite

\end{document}

